% Generated draft: requires mathematical and source review.
\documentclass[11pt]{article}
\usepackage[margin=25mm]{geometry}
\usepackage{fontspec}
\setmainfont{Latin Modern Roman}
\usepackage{amsmath,amssymb,mathtools,unicode-math}
\setmathfont{Latin Modern Math}
\usepackage{enumitem,needspace,xurl,hyperref}
\setlength{\parindent}{0pt}
\setlength{\parskip}{5pt}
\setlength{\emergencystretch}{3em}
\newcommand{\N}{\mathbb{N}}
\newcommand{\Z}{\mathbb{Z}}
\newcommand{\Q}{\mathbb{Q}}
\newcommand{\R}{\mathbb{R}}
\newcommand{\C}{\mathbb{C}}
\newcommand{\F}{\mathbb{F}}
\newcommand{\A}{\mathbb{A}}
\newcommand{\PP}{\mathbb{P}}
\newcommand{\E}{\mathbb{E}}
\newcommand{\eps}{\varepsilon}
\newcommand{\dd}{\,\mathrm{d}}
\begin{document}
\section*{UnsolvedMath ID: 30001342}
\textbf{Source problem code:} catalogue entry
\subsection{Definitions and mathematical statement}
The supplied record uses \(u(x,t)\) for an interface height, \(t\) for time, and \(F\) for a forcing strength. It does not specify the evolution equation, the spatial domain, the admissible forcing strengths, or the initial and boundary conditions.

For an interface evolution corresponding to forcing strength \(F\), the source describes pinning at a spatial point \(x\) by
\[
\lim_{t\to\infty}u(x,t)<+\infty,
\]
intended to mean that the limit exists and is finite, and describes absence of pinning at \(x\) by
\[
\lim_{t\to\infty}u(x,t)=+\infty.
\]
The question is to determine which forcing strengths satisfy the first condition and which satisfy the second.

These are the source's stated criteria, rather than a fully specified mathematical problem. In particular, the record does not state whether the conditions must hold at a fixed point, at every spatial point, or in another spatial sense. Although the background concerns random media, neither a probability model nor a probabilistic quantifier is supplied. Existence of the displayed limits and exhaustiveness of the two alternatives are not asserted by any supplied hypotheses. A precise classification therefore requires the missing evolution model and quantifiers.
\subsection{Short English statement}
Determine the forcing strengths for which the interface height converges to a finite limit and those for which it tends to \(+\infty\). The supplied record leaves the interface equation, spatial setting, initial and boundary data, and spatial and probabilistic quantifiers unspecified, so the classification problem remains incompletely formulated.
\subsection{Sources}
\begin{itemize}
\item Supplied local UnsolvedMath registry record, permanent ID \(30001342\), title \emph{Pinning Thresholds for Forced Interfaces}, source problem code \texttt{catalogue entry}, in the Oberwolfach Reports Open Problems collection. The record gives proposed year \(2009\) and names no proposer. Its label \texttt{corrected\_verified} is accompanied only by the explanation that mathematical delimiters were added without changing the content; it does not establish completeness of the formulation.
\item Supplied report citation: \emph{PDE and Materials (2010)}, Oberwolfach Reports, DOI \texttt{10.4171/owr/2009/42}, pp.~2335--2404. \url{https://doi.org/10.4171/owr/2009/42}. The background describes these reports as documenting mathematical workshops and their open-problem discussions. No problem-specific page or named author is supplied.
\item Supplied literature source, labelled ``pinning of interfaces in random media'': \url{https://arxiv.org/abs/0911.4254}.
\item Supplied literature source, labelled ``zero-mean random medium'': \url{https://doi.org/10.4171/IFB/455}.
\item Supplied literature source, labelled ``infinite pinning'': \url{https://arxiv.org/abs/2007.07605}.
\item Supplied dated curation, checked \texttt{2026-08-21}: the record labels the problem partially solved and reports that pinning and depinning have been proved in several dimensions and force regimes, including later zero-mean and infinite-pinning results, while the full phase diagram remains unsettled. This is the supplied literature triage derived from OWR curation, not an independently verified assessment. Its time-sensitive status and cited claims require independent verification.
\end{itemize}
\end{document}
